\section{Introduction}
%Neural network pruning has gained renewed interest with the success of deep neural networks in a variety of computer vision and machine learning tasks.
%Despite the success of deep neural networks in a variety of computer vision and machine learning tasks, they are highly overparametrized making them computationally expensive with excessive memory requirements.
Despite the success of deep neural networks in machine learning, they are often found to be highly overparametrized making them computationally expensive with excessive memory requirements.
%Specifically, modern neural networks are highly overparametrized (usually millions of learnable parameters) making them computationally expensive with huge memory requirements.
Pruning such large networks with minimal loss in performance is appealing for real-time applications, especially on resource-limited devices.
In addition, compressed neural networks utilize the model capacity efficiently, and this interpretation can be used to derive better generalization bounds for neural networks~(\cite{arora2018stronger}).
%which could encourage research on the generalization of neural networks~(\cite{arora2018stronger}).% and/or scalable incremental learning algorithms~(\cite{?}).\NOTE{refine this}

In network pruning, given a large reference neural network, the goal is to learn a much smaller subnetwork that mimics the performance of the reference network.
The majority of existing methods in the literature attempt to find a subset of weights from the pretrained reference network either based on a saliency criterion~(\cite{NIPS1988_119,lecun1990optimal,han2015learning}) or utilizing sparsity enforcing penalties~(\cite{chauvin1989back,Carreira-Perpiñán_2018_CVPR}).
%Unfortunately, since pruning is included as a part of the iterative optimization procedure, all these methods are prohibitively slow due to many expensive {\em prune -- retrain cycles}, and require highly heuristic design choices making them non-trivial to extend to new architectures and/or tasks.
Unfortunately, since pruning is included as a part of an iterative optimization procedure, all these methods require many expensive {\em prune -- retrain cycles} and heuristic design choices with additional hyperparameters, making them non-trivial to extend to new architectures and tasks.
 

%This unreasonably high training time hinders the use of pruning methods in large-scale applications.
%Thus, being able to prune the reference network once at the beginning and training a much smaller network with minimal loss in performance is a highly desirable goal.

In this work, we introduce a saliency criterion that identifies connections in the network that are important to the given task in a data-dependent way before training.
%Specifically, we discover important connections based on their influence on the loss function at a variance scaling initialization (we call this connection sensitivity).
Specifically, we discover important connections based on their influence on the loss function at a variance scaling initialization, which we call connection sensitivity.
%regardless of the weights of the connections. 
%Hence, our method can also be interpreted as an architecture selection method.
Given the desired sparsity level, redundant connections are pruned once prior to training (\ie, single-shot), and then the sparse pruned network is trained in the standard way.
%Our approach has several attractive properties as listed below:
Our approach has several attractive properties:
\begin{tight_itemize}
\vspace{-2mm}
\item \emph{Simplicity.}\;
    %Since the network is pruned once at random initialization, we have eliminated the need for pretraining and complex pruning schedules.
    Since the network is pruned once prior to training, there is no need for pretraining and complex pruning schedules.
    Our method has no additional hyperparameters and once pruned, training of the sparse network is performed in the standard way.
    % and in theory, the training can be accelerated using software libraries that support sparse matrix computations.
\item \emph{Versatility.}\;
	Since our saliency criterion chooses structurally important connections, it is robust to architecture variations.
	Therefore our method can be applied to various architectures including convolutional, residual and recurrent networks with no modifications. 
\item \emph{Interpretability.}\;
    %Our method chooses connections that are important to the given task where the task is defined by the loss function and the dataset.
	%By varying the data used for pruning and visualizing the retained connections, we verify that they are indeed important for the given task.
    Our method determines important connections with a mini-batch of data at single-shot. 
    %By varying this mini-batch used for pruning and visualizing the retained connections, we are able to verify that they are indeed important for the given task.
    By varying this mini-batch used for pruning, our method enables us to verify that the retained connections are indeed essential for the given task.   
\end{tight_itemize}

We evaluate our method on \mnist{}, \cifar{-10}, and \timagenet{} classification datasets with widely varying architectures.
Despite being the simplest, our method obtains extremely sparse networks with virtually the same accuracy as the existing baselines across all tested architectures.
%Furthermore, we provide interesting experiments by visualizing the important connections determined by our method, which, we believe sheds some light towards understanding neural networks.\NOTE{subject to results}
Furthermore, we investigate the relevance of the retained connections as well as the effect of the network initialization and the dataset on the saliency score. 



\SKIP{
Our approach has several attractive properties:
\begin{tight_itemize}
\vspace{-2mm}
\item \emph{Simplicity.}\;
    It requires no complex pruning -- learning schedule as well as any additional hyperparameters, being procedually simple and easy to use.
\item \emph{Promptitude.}\;
    It prune a given network only once, at random initialization without pretraining, achieving sparse network quickly.
\item \emph{Scalability.}\;
    It is architecture independent and model-agnostic and can be applied to various large-scale modern network types including convolutional, residual, recurrent networks without adjusting the algorithm.
\item \emph{Interpretability.}\;
    Our saliency measure is directly dependent on data, allowing one to interpret whether the right connections are survived from pruning.
\item \NOTE{Update, add, or remove. Once done, move around to fit into the intro.}
\end{tight_itemize}

General explanation
\begin{tight_itemize}
\item problem received new attention
\item big networks, computation speed, memory
\end{tight_itemize}

Problems and motivation
\begin{tight_itemize}
\item standard magnitude-based pruning: heuristic, requires fully-trained network
\item Improve the efficiency of pruning algorithms by introducing a new approach.
\item Bayesian can't be applied to large network, resnets, lstms.
\end{tight_itemize}

Approach
\begin{tight_itemize}
\item
\end{tight_itemize}

Contributions.
\begin{tight_itemize}
\item efficient pruning (\eg requires no pre-training or non-trivial pruning schedule)
\item new approach that can scale
\item demonstrate its applicability to various architectures (convolutional, fully-connected, recurrent, residual)
\item can be applied generally to training neural networks (sparse network encourages generalizability)
\end{tight_itemize}
}
